% Docking Protocol
\subsection{Docking Protocol for Energy Trend Analysis}

To assess whether molecules occupying candidate HEW sites show improved
binding energetics, we performed the following docking protocol:

\textbf{Step~1: MMFF94 Minimization.}
Each generated molecule was pre-minimized using the MMFF94 force field
(RDKit implementation) with 200 iterations and a force tolerance of
0.01~kcal/mol/\AA. Molecules failing MMFF94 parameterization were
minimized with UFF as a fallback.

\textbf{Step~2: Full Conformational Docking.}
Minimized molecules were converted to PDBQT format (OpenBabel, Gasteiger
charges) and docked into their respective protein pockets using AutoDock
Vina~1.2.3 with exhaustiveness~$=8$, generating up to 9 binding modes.
The docking box was centered on the pocket centroid with dimensions
$22$--$35$\,\AA\ per side (covering the pocket residues plus 8--10\,\AA\
padding).

\textbf{Step~3: Statistical Comparison.}
Molecules were grouped by whether they occupied at least one candidate HEW
site (compatible atom within 2.5\,\AA). The lowest Vina score for each
molecule was recorded. Group differences were assessed via the
Mann-Whitney $U$ test (two-sided) with Cliff's $\delta$ as the
effect size measure.

\textbf{Limitations.}
This protocol evaluates the \textit{generated} binding pose, not the
globally optimal pose after exhaustive conformational search.
Vina scores should be interpreted as estimates of relative binding
affinity within each pocket, not as absolute free energies.
